\chapter{Kéo Dài Giờ Ra Chơi}
\chapterintrobi{Chuông vào học đã dứt nhưng bước chân vẫn còn lưu luyến hành lang}{Stretch break time}{Có bạn coi giờ ra chơi là khái niệm co giãn tùy cảm xúc.}

\begin{lucbatbox}{Lục bát: Chuông xong vẫn đứng tám}
Chuông kia báo hết từ lâu\
Em còn đứng kể nhịp câu chưa dừng\
Hành lang lưu luyến vô cùng\
Lớp kia vào đủ, mình chung thủy ngoài\
Tới khi chạy vội mới hay\
Tiết sau bắt đầu chẳng đợi riêng ai
\end{lucbatbox}

\begin{tutuyetbox}{Thất ngôn tứ tuyệt: Luyến lưu ngoài cửa}
Giờ chơi em kéo thành hơi dài\
Chuông vào học cũng đứng nghe chơi\
Nếu tình yêu lớp sâu bằng thế ấy\
Chắc em vào sớm nhất trường thôi
\end{tutuyetbox}

\begin{ghichu}
Kéo dài giờ ra chơi rất vui trong lúc đó, nhưng thường làm phần đầu tiết sau nhập cuộc bằng một pha hớt hải.
\end{ghichu}

\begin{englishnote}
\textbf{break time}: giờ ra chơi.\par
\textbf{run back to class}: chạy về lớp.\par
\textbf{The bell rang, but we were still outside.}: Chuông đã reo mà bọn em vẫn còn ở ngoài.
\end{englishnote}